\documentclass[11pt]{article}
\usepackage[margin=.85in]{geometry}
\usepackage{amsmath,amssymb}
\usepackage[colorlinks=true,urlcolor=blue]{hyperref}
\setlength{\parindent}{0pt}
\setlength{\parskip}{.45em}
\setlength{\emergencystretch}{2em}
\newcommand{\LL}{\mathcal L}
\newcommand{\ch}{\chi_p}
\newcommand{\tw}{\widetilde\chi}
\newcommand{\gs}{\href{https://dms.umontreal.ca/~andrew/PDF/pretentious.pdf}{Granville--Soundararajan, JAMS 20 (2007), pp.\ 375--378}}
\newcommand{\lls}{\href{https://arxiv.org/pdf/1309.3595}{Lamzouri--Li--Soundararajan}}
\begin{document}
{\Large\bfseries The one-third endpoint for prime discriminant}\\[3pt]
{\large A test of the least-nonresidue bound on $p\equiv 7\pmod 8$}

\section*{1. The exact endpoint and the factor two}
Let $p\equiv23\pmod{24}$ be prime, $q=3p$, $\xi(n)=(n/3)$, $\ch(n)=(n/p)$, $\tw=\xi\ch$, and $S_p=\sum_{n\le p}\tw(n)$. Write $Y=(\log p)^2$, $\ell_2=\log\log p$ and $\LL=L(1,\ch)$. The Fourier expansion and class-number formula give
\begin{equation}
 S_p=\frac{\sqrt{3p}}{\pi}\sum_{n\ge1}\frac{\tw(n)\sin(2\pi n/3)}{n}
 =\frac{\sqrt p}{\pi}\LL=h(-p).
\end{equation}
Indeed $3p$ is a positive fundamental discriminant, so $\tw$ is its primitive even quadratic character and its Gauss sum is $+\sqrt{3p}$. The sine series is the Fourier expansion of the periodic partial sum, taken at $p/q=1/3$. Endpoint half-weights cause no discrepancy since $\tw(p)=0$. Now $\sin(2\pi n/3)=(\sqrt3/2)\xi(n)$ and $\ch(3)=1$. Summation over $3\nmid n$ gives $(1-\ch(3)/3)\LL=2\LL/3$. The last equality is Dirichlet's class-number formula for the fundamental discriminant $-p<-4$. Quadratic reciprocity identifies $\chi_{-p}=\ch$. \href{https://arxiv.org/pdf/1311.7556}{Bober--Goldmakher, Sections 1--2}. \lls, eq.\ (1.3).

In the proof of Granville--Soundararajan's Theorem 6 the nontrivial comparison character has conductor $m$. In the potentially extreme case here it is exactly $\xi$, with $m=3$. At the rational $b/r=1/3$ the only primitive odd character of conductor dividing $r$ is $\xi$. The principal character has the wrong parity. If $n_p>\eta Y$, the two characters $\tw,\xi$ agree at every prime $\le\log q$ other than $3$ eventually, and their common zero at $3$ contributes only a bounded term to the pretentious distance. The nonprincipal competitors are separated by the source's distance estimate. In the major-arc formula $N=q$ at this exact endpoint, $r/m=1$ and its product of correction factors is empty. The smooth sum is over $3\nmid n$, with character $\ch(n)$ there. The resulting upper majorant, to leading order, is
\[
 \frac{\sqrt{qm}}{\pi\varphi(r)}
       \prod_{v\le(\log q)^2,\ v\ne3}(1-1/v)^{-1}
 \sim\frac{3\sqrt p}{2\pi}\cdot\frac23\cdot e^\gamma\log((\log q)^2)
 \sim\frac{2e^\gamma}{\pi}\sqrt p\,\ell_2 ,
\]
where $v$ ranges over primes. The omitted factor at $3$ is exactly compensated. The ramified prime $p$ is already above the product cutoff. Lemma 6.4 discards larger prime factors using GRH in (6.7). Mertens' theorem on the remaining range gives $\log((\log q)^2)=2\log\log q$. No maximization over endpoints is needed for this obstruction. \gs, Lemmas 6.4--6.5 and Section 7.

On this very progression the known lower and conditional upper bounds give
\[
 \frac{e^\gamma}{\pi}\le
 \limsup_{p\equiv23\ (24)}\frac{S_p}{\sqrt p\,\ell_2}
 \le\frac{2e^\gamma}{\pi}.
\]
For the unconditional lower bound use Theorem 3.1 of Brink--Moree--Osburn with $\Delta=1,d=24,c=-1$. The two prescribed local factors both equal one since $c^*=1$. \href{https://arxiv.org/pdf/1003.1094}{Brink--Moree--Osburn, Theorem 3.1}.

\newpage
\section*{2. A class-number estimate from reduced forms}
Every class has a unique reduced positive definite form $[a,b,c']$ under the usual boundary convention. In particular $-a<b\le a$, $a\le c'$ and $a\le\sqrt{p/3}$. Let $R_p(a)$ count the residue classes $b\bmod 2a$ with $b^2\equiv-p\bmod 4a$. There is at most one reduced form per such class for fixed $a$. Conversely if $a<\sqrt p/2$, every root gives a reduced form by choosing $-a<b\le a$ and setting $c'=(b^2+p)/(4a)>a$. Primitivity is automatic for fundamental discriminant $-p$. Thus
\[
 \sum_{a\le\sqrt p/2}R_p(a)\ \le\ h(-p)\ \le\ \sum_{a\le\sqrt{p/3}}R_p(a).
\]
The endpoints are nonintegral. The local counts at a split prime power are $2$, and at an inert prime power they are zero. At the ramified prime $p$, the exponent-one count is $1$ and higher counts vanish. At $2$ there are two roots in the stated modulus convention at every positive exponent since $-p\equiv1\pmod8$. Equivalently $R_p(a)$ counts primitive integral ideals of norm $a$. Hence for $\Re s>1$,
\[
 \sum_{a\ge1}\frac{R_p(a)}{a^s}
       =\frac{\zeta(s)L(s,\ch)}{\zeta(2s)}.
\]
\href{https://magma.maths.usyd.edu.au/~watkins/papers/ALTMERGE15.pdf}{Watkins, primitive ideals and reduced forms}. \href{https://math.mit.edu/classes/18.785/2019fa/ProblemSet7.pdf}{MIT, reduction and uniqueness}.

Here is a direct estimate of this upper count, uniform in the prime and independent of GRH. For $X\asymp\sqrt p$,
\begin{equation}
 \sum_{a\le X}R_p(a)=\frac{X\LL}{\zeta(2)}
                       +O_\epsilon(p^{11/24+\epsilon}).
\end{equation}
Put $B(z)=\sum_{uv\le z}\ch(v)$, $C(t)=\sum_{v\le t}\ch(v)$. Burgess with $r=2$ gives $C(t)\ll_\epsilon t^{1/2}p^{3/16+\epsilon}$ (extend past $p$ by periodicity). Split the $v$ sum at $V$:
\[
 B(z)=\sum_{v\le V}\ch(v)\lfloor z/v\rfloor+
                \sum_{u\le z/V}(C(z/u)-C(V)).
\]
The first sum differs from $z\LL$ by $O_\epsilon(V+zp^{3/16+\epsilon}V^{-1/2})$ by partial summation of the tail. The second sum costs $O_\epsilon(zp^{3/16+\epsilon}V^{-1/2})$. Choose $V=(zp^{3/16})^{2/3}$ when $V\le z$. For smaller $z$, use $|B(z)|\ll z\log(2z)$ and $\LL\ll\log p$. Thus $B(z)-z\LL\ll_\epsilon p^{1/8+\epsilon}z^{2/3}$ for $1\le z\ll\sqrt p$. Sum $\mu(d)B(X/d^2)$ over $d\le\sqrt X$. The error converges since $\sum d^{-4/3}<\infty$, and extending $\sum\mu(d)/d^2$ costs $O(\sqrt X\log p)$. This proves (2). \href{https://doi.org/10.1112/plms/s3-13.1.524}{Burgess, On character sums and L-series II}.

At $X=\sqrt{p/3}$, inserting (1) into (2) gives the upper count
\[
 \frac{2\sqrt3}{\pi}\,h(-p)+O_\epsilon(p^{11/24+\epsilon}).
\]
The leading factor is about $1.103$. At $X=\sqrt p/2$ the lower count gives $(3/\pi)h(-p)$ with the same error. Thus these two genuinely arithmetic counts both scale with the same $L$-value. Applying the published explicit GRH upper bound for $h(-p)$ instead gives, for $p\ge10^{10}$,
\[
 \frac{\pi S_p}{e^\gamma\sqrt p}\le 2\ell_2-2\log2+1+2/\ell_2.
\]
The additive improvement is bounded. \lls, Theorem 1.5 and Corollary 1.3.

\newpage
\section*{3. How close an exceptional prime must come}
All estimates on this page assume Dirichlet GRH. Define the signed Euler-factor loss
\[
 D_p=\sum_{\substack{v\le Y\ {\rm prime}\\\ch(v)=-1}}
                         \log\frac{v+1}{v-1},\qquad
 \Psi_\chi(t)=\sum_{n\le t}\Lambda(n)\ch(n),\qquad
 E_p=\frac{\pi S_p}{e^\gamma\sqrt p}=e^{-\gamma}\LL .
\]
The short Euler product of Granville--Soundararajan gives $E_p=(1+o(1))\log Y\,e^{-D_p}$ uniformly. One can put this comparison on an absolute scale. For $x=Y$,
\begin{equation}
 E_p=e^{-D_p}\left(\log Y-\frac{\Psi_\chi(Y)}Y+\theta_p\right),
                    \qquad |\theta_p|\le1+o(1).
\end{equation}
In particular $|E_p-\log Y\,e^{-D_p}|\le2+o(1)$. \href{https://arxiv.org/pdf/math/0206031}{Granville--Soundararajan, Section 2, following Lemma 2.1}.

To prove (3), use $A_x=\sum_{n\le x}(\Lambda(n)\ch(n)/n)(1-n/x)$,
$a=(1-\ch(-1))/2$, $B=|\Re B(\ch)|$ in the notation of Lamzouri--Li--Soundararajan, and
$C_p=(\log(p/\pi)+\Gamma'/\Gamma((1+a)/2))/2$.
Their Lemma 2.3 gives
\[
 |B-C_p+A_x|\le 2B/\sqrt x+O((B+\log p+\log x)/x).
\]
Indeed their parity factor is $\frac12(\Gamma'/\Gamma)((1+a)/2)+O((1+\log x)/x)$. At $x=Y$ the same lemma and $A_x=O(\log x)$ give $B=\frac12\log p+O(\log Y)$. Their Lemma 2.5 gives
\[
 \log\LL=\sum_{n\le x}\frac{\Lambda(n)\ch(n)}{n\log n}
                 \frac{\log(x/n)}{\log x}
       +\frac{C_p-B}{\log x}
       +O\!\left(\frac{B}{\sqrt x\log^2 x}+\frac1{x\log^2 x}\right),
\]
with sums containing $\log n$ starting at $n=2$. Combining the two estimates gives
\[
 \left|\log\LL-T_Y(\ch)\right|\le\frac{1+o(1)}{\log Y},
 \quad T_x(\ch)=\sum_{n\le x}\Lambda(n)\ch(n)
                  \left(\frac1{n\log n}-\frac1{x\log x}\right).
\]
The omitted higher powers in passing from the first part of $T_Y$ to the logarithm of the short Euler product cost $O(Y^{-1/2})$ by splitting the prime bases at $\sqrt Y$ and summing their geometric tails. Mertens' product with the PNT remainder is $e^\gamma(\log Y+o(1))$. Since $|\Psi_\chi(Y)|/Y\le1+o(1)$, exponentiation gives (3). \lls, Lemmas 2.3 and 2.5.

Suppose $n_p>\eta Y$ for fixed $0<\eta\le1$. Every weight in $T_Y$ is nonnegative, so its smallest allowed value is bounded below by putting minus signs on every term above $\eta Y$. The ordinary PNT gives
\begin{align*}
 T_Y(1)&=\log\log Y+\gamma-(1+o(1))/\log Y,\\
 \sum_{\eta Y<n\le Y}\Lambda(n)\left(\frac1{n\log n}-\frac1{Y\log Y}\right)
   &=\big(\log(1/\eta)-1+\eta+o_\eta(1)\big)/\log Y .
\end{align*}
Subtract twice the latter quantity, use the lower estimate for $\log\LL$, and exponentiate to obtain
\begin{equation}
 n_p>\eta Y\quad\Longrightarrow\quad
 E_p\ge 2\ell_2-2\log(1/\eta)-2\eta+o_\eta(1).
\end{equation}
Thus $2\ell_2-E_p\to+\infty$ on this progression would suffice for the original little-$o$ bound there. Equation (4) supplies the implication. The uniform divergence is open in this calculation.

\newpage
\section*{4. The size of a fixed saving}
In the original units a uniform relative saving $\delta\in(0,1)$ means
\[
 \limsup_{p\equiv23\ (24)}\frac{S_p}{(2e^\gamma/\pi)\sqrt p\,\ell_2}
       \le1-\delta.
\]
The short Euler product shows this is equivalent to $\liminf D_p\ge-\log(1-\delta)$. It would in fact imply
\[
 \limsup_{p\equiv23\ (24)}
        \frac{\log n_p}{\log\log p}\le2\sqrt{1-\delta}.
\]
Indeed if all primes up to $(\log p)^\alpha$ are residues, for fixed $0<\alpha<2$, then $D_p\le2\log(2/\alpha)+o(1)$ by the prime harmonic sum and $\log((v+1)/(v-1))=2/v+O(v^{-3})$. Any fixed $2\sqrt{1-\delta}<\alpha<2$ is excluded eventually. In comparison, (4) only needs a divergent additive saving. By (3) this weaker endpoint saving is equivalent to $(\log Y)D_p\to\infty$ since $1-e^{-d}\asymp\min(d,1)$ for $d\ge0$.

\textbf{Check on the exponential weight.} Let $X=0.9Y$, $x=tY$, $t=0.9/2.452$,
and $\Psi_-(X)=\sum_{n\le X,\,\ch(n)=-1}\Lambda(n)$. Apply Lemma 6.1 of Lamzouri--Li--Soundararajan with $K(s)=\Gamma(s+1/2)$, whose inverse Mellin transform is $\sqrt u e^{-u}$. For the two characters $1$ (principal mod $p$) and $\ch$, the upper bound for the combined sum is $\sqrt x+(C+o(1))\log p$ in the normalization $x^{-1/2}\sum_n\Lambda(n)(\chi_0(n)+\ch(n))e^{-n/x}$. Here
\[
 C=(2\pi)^{-1}\int_{\mathbb R}|\Gamma(1/2+is)|\,ds
   =\frac{\Gamma(1/4)}{2\pi\Gamma(3/4)}.
\]
All combined coefficients are nonnegative. Discard $n>X$ and apply the PNT to the unsigned sum up to $X$. The retained negative terms have $e^{-n/x}\le1$. This yields
\[
 \liminf_p\frac{\Psi_-(0.9Y)}{Y}\ge
 \frac{t(1-2e^{-0.9/t})-C\sqrt t}{2}
 =0.0092707\ldots .
\]
The integral from $0.9Y$ to $Y$ has lower limit at least $0.0092707\ldots/9>0.001$ after weighting by $t'^{-2}dt'$. Prime powers of exponent at least two contribute $o(Y)$ to the displayed prime-power count. All these statements hold for all odd primes. For any fixed $0<\eta<0.9$, even assigning negative sign to every prime in $(\eta Y,0.9Y]$ accounts for only $O_\eta(1/\log Y)$ of $D_p$. Thus information confined to such an interval can supply at most a bounded additive saving in (3). \lls, Lemma 6.1 and Section 6.3.

\textbf{Scope.} For $p\equiv7\pmod{24}$ one has $(3/p)=-1$. These and $p\equiv23\pmod{24}$ exhaust the primes $p\equiv7\pmod8$. A proof of the remaining uniform additive saving on $23\pmod{24}$ would therefore settle the whole $7\pmod8$ branch. The estimates here give no improvement to $n_p$ on $23\pmod{24}$ and give no bound beyond the existing ones on $1\pmod8$. The deductions about the signed Euler product and reduced-form count are hand derivations.
\end{document}
